\begin{frame}{1969 --- 첫 겨울의 시작}
  \begin{itemize}
    \item 민스키·페퍼트가 단층 퍼셉트론은
      \textbf{XOR}(배타적 논리합, ``둘 중 정확히 하나만 참일 때만 참'')조차
      표현할 수 없음을 \textbf{수학적으로 증명}.
    \item 결과: 연구비가 끊기고 연구자 수가 급감 ---
      \kb{``첫 AI 겨울''}이 10년 넘게 이어짐.
  \end{itemize}
  \vskip0.3em
  {\small 책의 결론부는 ``아직 알고리즘이 부족하다''가 아니라 다층으로의 확장 자체를 막다른 길로 본다:}\par
  \vskip0.3em
  \hspace*{0.5em}{\color{ksagolddk}\vrule width 3pt}\hspace{0.6em}%
  \begin{minipage}[c]{0.88\linewidth}
    \textit{\ldots we consider it to be an important research problem to elucidate
    (or reject) our intuitive judgement that the extension is \textbf{sterile}.}\par
    \vskip0.2em
    {\small\color{ksagray}``\ldots 이 확장이 \textbf{쓸모없다(sterile)}는 우리의 직관적 판단을 규명하는
    (또는 반박하는) 것이 중요한 연구 과제라고 본다.''}\par
    \vskip0.2em
    \hfill{\small\color{ksagray}--- Minsky \& Papert, \textit{Perceptrons} (1969), p.~232}\par
  \end{minipage}\par
  \vskip0.2em
  \begin{block}{흥미로운 점}
    답(층을 여러 개 쌓으면 XOR도 표현 가능)은 \textbf{이미 알려진 수학} 안에
    있었지만, ``그 여러 층을 \textbf{어떻게 학습}시킬 것인가''라는
    \textbf{계산적 문제}가 안 풀려서 겨울이 그토록 길어졌다.
  \end{block}
  \vskip0.4em
  {\tiny\color{ksagray} Minsky, M. \& Papert, S. (1969). \textit{Perceptrons:
  An Introduction to Computational Geometry}. MIT Press.}
\end{frame}
